\section{Discussion: scope, limitations, and next steps}
\label{sec:discussion}

\subsection{What this draft establishes}
This paper contributes a mathematics instantiation of the \SBT dictionary of \cite{SixBirdsFramework} with three deliverables:
\begin{itemize}
  \item a \emph{formal vocabulary} for mathematics-as-closure (Section~\ref{sec:dictionary} and Section~\ref{sec:packaging-engine}),
  \item \emph{machine-checked anchors} (Lean) and falsification-first diagnostics (Python) that can fail,
  \item an \emph{artifact contract} that ties reported numbers in the PDF to snapshot-visible pointers in the repository (Section~\ref{sec:reproducibility}).
\end{itemize}
The exhibits are intended to be small but sharp: they illustrate how SBT-style feasibility pressure (staging + ledger + packaging + constraints) can be expressed in ordinary mathematical workflows.

\subsection{Limitations and guardrails}
Several limitations are deliberate and should be understood as guardrails rather than omissions.

\paragraph{Diagnostics are not theorems.}
The Python experiments are designed as feasibility stress tests, not as proofs of uniqueness.
For example, the stencil-selection results support the claim that adding a Leibniz-defect gate selects derivative-like operators under the tested conditions, but they do not rule out exotic behaviors outside those conditions.

\paragraph{Prime closure results are baseline, not analytic continuation.}
Exhibit~B-1 separates a convergence control regime from the critical strip and reports catastrophic mismatch for naive staging in the strip.
This is consistent with classical domain-of-convergence facts; it does not establish new results about $\zeta$ and should not be read as evidence for conjectures about its zeros.

\paragraph{Positivity toy is a motif demonstration.}
The passivity toy illustrates a pattern---tightening a positivity constraint can confine zeros to a symmetry locus---but it does not claim generality.
Its role is motivational: it shows why ``zero confinement to a self-dual manifold'' naturally invites positivity/passivity ledgers.

\paragraph{SBT is used as a dictionary, not a replacement foundation.}
This paper does not propose new axioms or attempt to subsume all of mathematics under one formalism.
Instead it makes explicit a set of closure roles that are already present (often implicitly) in mathematical practice.

\subsection{Upgrade points}
The exhibits suggest concrete upgrade paths if one wishes to strengthen claims.

\paragraph{Stronger analytical closure claims.}
For calculus, one can replace heuristic feasibility gates with classical stability/consistency hypotheses from numerical analysis, and explicitly state the regularity classes under which the ledger terms vanish.
For protocol mismatch, one can connect the observed scaling to known truncation-order estimates.

\paragraph{Richer prime closure ledgers.}
Exhibit~B-1 uses deliberately simple staging and packaging.
For $\zeta(s)$ in the critical strip, a natural upgrade is to use the approximate functional equation as an additive staged protocol (valid in the strip) and to compare against smoothed prime-sum / log-derivative ledgers in regimes where such representations are controlled \cite{Titchmarsh1986}.
More informative ledgers could compare additional invariants (e.g.\ log-derivative statistics), incorporate smoothing optimized for convergence acceleration, or restrict to regimes where the protocols admit controlled continuation.
The present draft should be treated as a baseline diagnostic that clarifies what fails without additional structure.

\paragraph{Formalization interface improvements.}
The Lean anchors are intentionally minimal.
Future versions can add small ``glue lemmas'' that align the informal narrative even more tightly with the formal objects (without attempting to formalize all of analysis in a proof assistant).

\subsection{Why the closure lens is still useful}
Even at this modest level, the closure lens supplies two pragmatic benefits.

First, it separates \emph{definitions} from \emph{feasibility}.
A packaged object is not merely a symbolic completion; it is a completion that is stable under refinement and compatible with chosen constraints.

Second, it encourages controlled comparisons.
The same route-mismatch and defect-ledger ideas can be applied to other mathematical packaging problems (measure and integration, functional analysis limits, geometric patching, spectral constructions), where it is otherwise easy to conflate ``a formula exists'' with ``a stable closure exists.''

\subsection{Closing}
The slogan ``To Count a Stone with Six Birds'' is literal in this instantiation:
counting is a staged discrete protocol; calculus is the stable compression of that protocol under an accounting ledger; and higher layers arise by repeating packaging with new constraints and new mismatch diagnostics.
The intent of the paper is to make those roles explicit, auditable, and reusable.
